(1) 令上层 \(\phi_1=A_1e^{-ky}e^{i(kx-\omega t)}\)，下层 \(\phi_2=A_2e^{ky}e^{i(kx-\omega t)}\)，均满足Laplace方程并在远离界面处衰减。线性运动学条件 \((\partial_t+U_j\partial_x)\eta=\partial_y\phi_j|_0\) 给 \(-i(\omega-kU_1)\eta_0=-kA_1\) 与 \(-i(\omega-kU_2)\eta_0=kA_2\)，因此 \(A_1=i(\omega-kU_1)\eta_0/k\)、\(A_2=-i(\omega-kU_2)\eta_0/k\)。正负指数不能颠倒，否则波在无穷远发散。

(2) 线性Bernoulli给界面两侧 \(p_j'=-\rho_j[(\partial_t+U_j\partial_x)\phi_j+g\eta]\)。按曲率取压差 \(p_2'-p_1'=\sigma k^2\eta\)，代入第一问系数并整理得 \(\rho_1(\omega-kU_1)^2+\rho_2(\omega-kU_2)^2=(\rho_2-\rho_1)gk+\sigma k^3\)。定义 \(\bar U=(\rho_1U_1+\rho_2U_2)/(\rho_1+\rho_2)\)，完成平方得题示色散式。虚部来自根号内负值而非对流项，故仅靠多普勒平移不会产生不稳定。

(3) 根号内非负时两个实频率稳定，失稳条件为 \(\frac{\rho_1\rho_2}{\rho_1+\rho_2}(\Delta U)^2>(\rho_2-\rho_1)g/k+\sigma k\)。对正 \(k\) 最小化右侧：\(k_*=\sqrt{(\rho_2-\rho_1)g/\sigma}\)，最小值 \(2\sqrt{(\rho_2-\rho_1)g\sigma}\)。所以 \(\boxed{(\Delta U_{\rm cr})^2=\frac{2(\rho_1+\rho_2)}{\rho_1\rho_2}\sqrt{(\rho_2-\rho_1)g\sigma}}\)。该阈值依赖稳定密度分层与非零表面张力，不能把Rayleigh–Taylor无剪切条件直接替换为此结果。

(4) 若 \(\Delta U=0\)，根号内为正，得到正常重力—毛细波 \(\omega=kU\pm\sqrt{[(\rho_2-\rho_1)gk+\sigma k^3]/(\rho_1+\rho_2)}\)。若 \(\sigma=0\)，任意非零剪切总有足够大的 \(k\) 使 \(k^2\) 负项占优。若 \(g=\sigma=0\)，增长率为 \(\gamma=k|\Delta U|\sqrt{\rho_1\rho_2}/(\rho_1+\rho_2)\)，随 \(k\) 无界增长。现实中有限剪切层厚度、黏性、可压缩性与表面张力等微观机制使极短波理想化失效；该发散不代表可实现无限增长率。
